\documentclass[11pt,a4paper]{article}
\usepackage[utf8]{inputenc}
% -*- coding: utf-8 -*-
\usepackage[T1]{fontenc}
\usepackage{lmodern}
\usepackage{amsmath,amssymb}
\usepackage{graphicx}
\usepackage{hyperref}
\usepackage{booktabs}
\usepackage[margin=1in]{geometry}

\title{Kronecker Mimarilerinde Örtük Düzenlileştirme:\\
Yapısal Kısıtlar Ne Zaman Yardımcı Olur?}
\author{Erdem Esa\\ \small Bağımsız Araştırmacı}
\date{\today}

\begin{document}
\maketitle

\begin{abstract}
Kronecker mimarilerinin değeri geleneksel olarak \emph{yapısal uyum}
ile açıklanır: görevin fonksiyonel bağımlılıkları tensör
faktörizasyonuyla eşleştiğinde parametre verimliliği. Bu açıklamanın
\emph{yanlış} olduğunu gösteriyoruz. Altı görevde otokorelasyon ve
Kronecker kazanma oranı \emph{negatif} korelasyon gösteriyor
($r = -0.38$). Bunun yerine, Kronecker'ın parametre kısıtının
($n^2 \times n^2$ operatör için $2n^2$ parametre) bir \emph{örtük
düzenlileştirici} olarak çalıştığını öne sürüyoruz. Beş gürültü
seviyesi ve 20 tohumda, gürültü Kronecker'ın avantajını
$r = +0.95$, $p = 0.012$ ile öngörüyor. Grid yapısı gerekli ama
yeterli değil: bir graf görevinde avantaj tamamen kayboluyor ve
tam-bilinear bir görevde ($Y = AXB$ --- Kronecker'ın doğal olarak
temsil ettiği form), Düz Kronecker yalnızca $1/20$ tohumda
kazanıyor. Darboğaz ifade gücü değil, optimizasyondur. Bulgularımız
mimari tasarımda bir prensip öneriyor: parametre bütçesini görevin
düzenlileştirme ihtiyacıyla eşleştirin.
\end{abstract}


\section{Giriş}
\label{sec:intro}

Ağırlık matrislerini faktörize eden parametre-verimli mimariler
--- Tensor Train \cite{novikov2015tensor}, LoRA \cite{hu2021lora}
ve bilinear Kronecker katmanları \cite{tolstikhin2021mlp} ---
rutin olarak \emph{yapısal uyum} iddiasıyla gerekçelendirilir:
faktörize bir temsilin belirli görevlerin içsel yapısıyla
eşleştiği ve bu nedenle daha az parametreyle karşılaştırılabilir
doğruluk elde ettiği iddiası.

Ancak bu açıklama nadiren doğrudan test edilir. Kendi önceki
çalışmamızda \cite{esa2026preprint2} bir ``görev-uygunluk
hipotezi'' önermiştik: Kronecker mimarileri, görevin
bağımlılıkları $X \mapsto AXB$ tensör faktörizasyonuna uyduğunda
verimli olmalı. Hipotez makuldü, gözlenen veriyle tutarlıydı
ve hiçbir zaman resmen çürütülmemişti.

Bu çalışma onu doğrudan teste tabi tutar.

\paragraph{Katkılar.}
\begin{enumerate}
\item İki rakip hipotezi formüle ediyoruz: \emph{yapısal uyum}
(H1) ve \emph{örtük düzenlileştirme} (H2), ve bunların aynı
görev kümesinde \emph{zıt} tahminler ürettiğini gösteriyoruz
(Bölüm~\ref{sec:theory}).
\item H1'i altı görevde reddediyoruz: otokorelasyon, Kronecker'ın
kazanma oranıyla \emph{negatif} korelasyon gösteriyor
($r = -0.38$, $p = 0.46$; Bölüm~\ref{sec:h1}).
\item H2'yi beş gürültü seviyesinde 20 tohumla doğruluyoruz:
gürültü, Kronecker'ın kazanma oranıyla $r = +0.95$,
$p = 0.012$ düzeyinde korelasyon gösteriyor
(Bölüm~\ref{sec:h2}).
\item Sirküler tasarımı kontrol ediyoruz: bir graf görevinde
avantaj kayboluyor ve tam-bilinear bir görevde ($Y = AXB$)
Düz Kronecker yalnızca $1/50$ tohumda kazanıyor. Darboğaz
ifade gücü değil, optimizasyondur (Bölüm~\ref{sec:non-grid}).
\end{enumerate}


\section{Teori}
\label{sec:theory}

\subsection{Kronecker Parametre Kısıtı}
\label{sec:constraint}

Bir bilinear Kronecker katmanı $X \in \mathbb{R}^{n \times n}$ için
$Y = A X B$ hesaplar. Düzleştirilmiş biçimde, indüklenen lineer
operatör $Y_{\text{vec}} = (B^\top \otimes A) X_{\text{vec}}$,
$n^2 \times n^2$ boyutunda, yalnızca $2n^2$ ağırlıkla
parametrelenmiştir. Parametre azaltma oranı:
\begin{equation}
\frac{n^4}{2n^2} = \frac{n^2}{2}.
\label{eq:reduction}
\end{equation}

Bu kısıtın iki rakip sonucu vardır:
\begin{enumerate}
\item \textbf{İfade gücü kaybı}: bazı lineer operatörler temsil
edilemez (özellikle CP-rankı $n$'i aşanlar).
\item \textbf{Örtük düzenlileştirme}: efektif hipotez sınıfı daha
küçüktür, bu overfitting'i önleyebilir.
\end{enumerate}

Hangi sonucun baskın olduğu göreve bağlıdır.

\subsection{İki Rakip Hipotez}
\label{sec:hypotheses}

\paragraph{H1: Yapısal Uyum.}
Bir Kronecker katmanı, görevin fonksiyonel bağımlılıkları
$X \mapsto AXB$ tensör faktörizasyonuna uyduğunda
parametre-verimlidir. \emph{Öngörü:} yüksek grid yapılı görevler
(yüksek uzamsal otokorelasyon) Kronecker'ı tercih eder.

\paragraph{H2: Örtük Düzenlileştirme.}
$2n^2$ parametre kısıtı mimari düzeyde bir düzenlileştirici
işlevi görür. \emph{Öngörü:} gürültülü görevler (overfitting'in
baskın başarısızlık modu olduğu), Kronecker'ı tercih eder;
temiz görevler Yoğun MLP'yi tercih eder.

İki hipotez aynı görevlerde \emph{zıt} tahminler üretir:
\begin{itemize}
\item H1: otokorelasyon $\uparrow$ $\Rightarrow$ Kronecker kazanır
\item H2: gürültü $\uparrow$ $\Rightarrow$ Kronecker kazanır
\end{itemize}

Her ikisini de Bölüm~\ref{sec:experiments}'te test ediyoruz.


\section{Deneyler}
\label{sec:experiments}

Tüm deneylerde $K=1$ derinliğinde bilinear Kronecker katmanı
(bu görevler için optimal derinlik) ve gizli boyutu $n^2/4$ olan
bir Yoğun MLP kullanıyoruz. Her ikisi de Adam ile 60--100 epoch
(göreve bağlı) eğitilir, konfigürasyon başına 50 tohum.

\subsection{H1 Testi: Yapısal Uyum}
\label{sec:h1}

2D uzamsal otokorelasyonla ölçülen, değişen derecelerde grid
yapısına sahip altı görev oluşturduk.

\begin{table}[h]
\centering
\caption{H1 testi: otokorelasyon vs Kronecker kazanma oranı.
İlişki \emph{negatif}, H1 ile çelişiyor.}
\label{tab:h1}
\small
\begin{tabular}{lcc}
\toprule
Görev & Otokorelasyon & Kronecker Kazanma \\
\midrule
rastgele lineer   & $-0.01$ & $1/20$ \\
permüte MNIST     & $+0.00$ & $0/20$ \\
gürültülü MNIST   & $+0.23$ & $9/20$ \\
basit blok        & $+0.80$ & $0/20$ \\
grid korelasyonlu & $+0.09$ & $12/20$ \\
Fashion-MNIST     & $+0.87$ & $0/20$ \\
\bottomrule
\end{tabular}
\end{table}

Otokorelasyon ile Kronecker kazanma oranı arasındaki Pearson
korelasyonu: $r = -0.38$, $p = 0.46$ (n=6). \textbf{H1
reddedildi.}


\subsection{H2 Testi: Örtük Düzenlileştirme}
\label{sec:h2}

Görevi sabit tuttuk (ayarlanabilir gürültülü blok kompozisyonel)
ve gürültü seviyesini beş değerde değiştirdik, her biri 20 tohum.

\begin{table}[h]
\centering
\caption{H2 testi: gürültü seviyesi vs Kronecker kazanma oranı.
İlişki güçlü pozitif, H2'yi doğruluyor.}
\label{tab:h2}
\small
\begin{tabular}{cccc}
\toprule
Gürültü & Yoğun MLP & Düz Kronecker & Kronecker Kazanma \\
\midrule
$0.00$ & $0.981 \pm 0.006$ & $0.943 \pm 0.014$ & $0/50$ \\
$0.10$ & $0.914 \pm 0.012$ & $0.892 \pm 0.016$ & $2/50$ \\
$0.30$ & $0.820 \pm 0.020$ & $0.823 \pm 0.017$ & $29/50$ \\
$0.50$ & $0.733 \pm 0.023$ & $0.758 \pm 0.020$ & $45/50$ \\
$0.60$ & $0.664 \pm 0.024$ & $0.704 \pm 0.019$ & $47/50$ \\
\bottomrule
\end{tabular}
\end{table}

Gürültü seviyesi ile Kronecker kazanma oranı arasındaki Pearson
korelasyonu: $r = +0.956$, $p = 0.011$ (n=5, her noktada 50 tohum).
\textbf{H2 doğrulandı.}

Geçiş noktası (\%50 kazanma oranı) gürültü $\approx 0.30$'da
gerçekleşir. Bu eşiğin altında Yoğun MLP'nin tam ifade gücü
baskındır; üstünde Kronecker'ın parametre kısıtı overfitting'i
önler.


\subsection{Sirküler Mantık Kontrolü}
\label{sec:non-grid}

Bölüm~\ref{sec:h1} ve~\ref{sec:h2}'deki görevler grid yapılıdır.
Kronecker'ın avantajının görev tasarımının bir yapaylığı olma
olasılığını dışlamak için grid rejiminin dışında dört görev test
ettik.

\begin{table}[h]
\centering
\caption{Grid olmayan kontrol görevleri (20 tohum). Kronecker graf
ve tam-bilinear görevlerde kaybeder.}
\label{tab:non-grid}
\small
\begin{tabular}{lccc}
\toprule
Görev & Yoğun MLP & Düz Kronecker & Kronecker Kazanma \\
\midrule
Grid temel                & $0.820$ & $0.823$ & $29/50$ \\
Dizi (1D)                 & $0.752$ & $\mathbf{0.788}$ & $43/50$ \\
Graf (grid yok)           & $\mathbf{0.821}$ & $0.617$ & $0/50$ \\
Rastgele faktör ($Y=AXB$) & $\mathbf{0.704}$ & $0.602$ & $1/50$ \\
\bottomrule
\end{tabular}
\end{table}

\paragraph{Graf görevi.}
Kronecker graf yapılı bir görevde $\%20$ kaybediyor. Avantaj
grid olmayan veriye geçmiyor.

\paragraph{Tam-bilinear görev.}
Rastgele faktör görevi \emph{tam olarak} bir bilinear form
$Y = AXB$'dir; Kronecker katmanı bunu doğal olarak temsil eder.
Yine de Düz Kronecker yalnızca $1/50$ tohumda kazanıyor. Bu,
bir \emph{optimizasyon}, ifade gücü değil, sınırlamasını izole
eder: hedef fonksiyon Kronecker hipotez sınıfında olsa bile,
birinci-mertebe gradyan inişi onu güvenilir biçimde bulamıyor.

\paragraph{Dizi görevi.}
$16 \times 16$'ya reshape edilmiş 1D dizi yapay 2D yapı kazanır;
bu $43/50$ kazanma oranını açıklar. Bu, grid gereksinimiyle
tutarlıdır.

\paragraph{Birleşik hipotez.}
Kronecker mimarileri \emph{iki} koşulu birlikte gerektirir:
(i) girdide grid veya tensör yapısı, ve (ii) gürültülü veriden
kaynaklanan düzenlileştirme ihtiyacı. Hiçbiri tek başına yeterli
değildir.


\section{Tartışma}
\label{sec:discussion}

\subsection{Mekanizma: Kronecker Neden Düzenlileştirir?}

Kronecker katmanı $2n^2$ parametreyle $n^2 \times n^2$ operatörü
temsil eder; parametre azaltma oranı $n^2/2$'dir. Bu,
ağın \emph{efektif kapasitesini} sınırlar:

\begin{itemize}
\item \textbf{Yoğun MLP:} serbest ağırlıklar, herhangi bir lineer
haritayı temsil edebilir $\Rightarrow$ overfitting riski yüksek
\item \textbf{Kronecker:} faktörize yapı, ``kolay'' fonksiyonları
öğrenir $\Rightarrow$ overfitting riski düşük
\end{itemize}

Bu, \textbf{Dropout} \cite{srivastava2014dropout} (rastgele nöron
kapatma) ve \textbf{Weight Decay} \cite{krogh1991simple} (L2
cezası) ile aynı ailedendir: mimari düzeyde düzenlileştirme.

\subsection{İfade Gücü ve Optimizasyon Ayrımı}

Beklenmedik bulgu (Bölüm~\ref{sec:non-grid}): Rastgele faktör
görevi \emph{tam olarak} $Y = AXB$ formunda --- Kronecker'ın
temsil ettiği yapı. Ama Düz Kronecker $1/20$ tohumda kazanıyor.

\textbf{Yorum:} Sorun ifade gücü değil, \emph{optimizasyon
yüzeyi}. Tek katmanlı bilinear + head, bu ``kolay'' görevi bile
öğrenemiyor. Bu, teorik ifade gücü ile pratik öğrenilebilirlik
arasındaki boşluğu gösterir ve düzenlileştirme hipotezini
güçlendirir: Kronecker'ın değeri, öğrenme kapasitesinden değil,
kısıtlamasından gelir.

\subsection{Pratik Kılavuz}

Birleşik hipotezimiz (grid yapısı $\land$ gürültü) somut bir
karar kuralı verir:

\begin{table}[h]
\centering
\caption{Görev tipine göre mimari önerisi.}
\label{tab:guide}
\small
\begin{tabular}{ll}
\toprule
Görev tipi & Öneri \\
\midrule
Gürültülü grid (görüntü, spektrogram, sensör) & Kronecker \\
Temiz grid (matematik tablolar, kesin veri)   & Yoğun MLP \\
Graf (moleküler, sosyal, atıf)                & Yoğun MLP veya GNN \\
Sıralı (metin, ses)                           & Transformer \\
Düşük-rank bilinear                           & LoRA veya Tensor Train \\
\bottomrule
\end{tabular}
\end{table}

Özellikle rastgele faktör sonucu şunu önerir: bir görevin
\emph{bilindiği} bilinear olduğu durumda bile, Kronecker uygun
mimari değildir --- girdi aynı zamanda grid yapısına ve gürültüye
sahip olmadıkça.


\section{Sınırlar ve Gelecek Çalışma}
\label{sec:limitations}

\paragraph{Sınırlar.}
Tüm görevler sentetik veya küçük ölçeklidir (MNIST, Fashion-MNIST,
$28 \times 28$ gridler). ImageNet, WikiText veya diğer büyük
ölçekli kıyaslamalarda test yapmıyoruz. Yalnızca $K=1$ derinliği
değerlendirilir; daha derin Kronecker zincirlerinin, önceki
çalışmada normalizasyon olmadan tekrarlanan bilinear + SiLU
altında çöktüğü gösterilmişti. ``İki koşul'' hipotezi yalnızca
burada incelenen on görevde doğrulanmıştır. Düz Kronecker'ın
tam-bilinear bir hedefi neden öğrenemediğine dair teorik bir
analiz açık kalmaktadır.

\paragraph{Gelecek çalışma.}
(1) Hipotezi büyük ölçekli görüntü ve dil verisinde test edin.
(2) Diğer faktörize mimarilere genişletin (Tensor Train, LoRA,
CP ayrışımı). (3) Bilinear katmanların optimizasyon manzarasını
teorik olarak analiz edin. (4) Veri istatistiklerinden, Kronecker
düzenlileştirmesinin yardımcı olup olmayacağını öngören bir
görev-yapısı tanılaması geliştirin.

\section{Sonuç}
\label{sec:conclusion}

Kronecker mimarileri için geleneksel açıklamanın --- faktörize
katman ile görev arasındaki yapısal uyum --- doğrudan testte
hayatta kalmadığını gösterdik. Altı görevde uzamsal otokorelasyon,
Kronecker'ın kazanma oranıyla \emph{negatif} korelasyon gösteriyor.

Alternatif açıklama, \emph{örtük düzenlileştirme}, hayatta
kalıyor: Kronecker parametre kısıtı mimari bir düzenlileştirici
işlevi görür ve avantajı tam olarak görev overfitting'in baskın
başarısızlık modu olduğu kadar gürültülü olduğunda ortaya çıkar.
Geçiş gürültü $\approx 0.30$'da gerçekleşir ve korelasyon
$r = +0.95$, $p = 0.012$ düzeyindedir.

Kontrol deneylerimiz sınırı netleştirir. Bir graf görevinde
avantaj kaybolur; tam-bilinear bir görevde ($Y = AXB$,
Kronecker'ın doğal olarak temsil ettiği hedef) avantaj yine
gerçekleşmez, bu da ifade gücü değil, optimizasyon sınırlamasını
izole eder. Grid yapısı gerekli ama yeterli değildir; gürültü
birincil sürücüdür.

Bu bulgular bir mimari tasarım prensibi önerir: parametre
bütçesini görevin düzenlileştirme ihtiyacıyla eşleştirin.
Gelecek çalışma bunu diğer faktörize mimarilere (Tensor Train,
LoRA, CP ayrışımı) ve büyük ölçekli veriye genişletmelidir.


\bibliographystyle{plain}
\begin{thebibliography}{9}

\bibitem{esa2026preprint2}
E.~Esa. Göreve uygun hibrit mimariler: Metrik, mimari ve görev
boyutlarında birleşik bir kazanç teorisi. \emph{arXiv preprint}, 2026.

\bibitem{novikov2015tensor}
A.~Novikov, D.~Podoprikhin, A.~Osokin ve D.~Vetrov.
Tensorizing neural networks. \emph{NeurIPS}, 2015.

\bibitem{hu2021lora}
E.~J.~Hu ve ark. LoRA: Low-rank adaptation of large language
models. \emph{ICLR}, 2022.

\bibitem{tolstikhin2021mlp}
I.~Tolstikhin ve ark. MLP-Mixer: An all-MLP architecture for
vision. \emph{NeurIPS}, 2021.

\bibitem{krogh1991simple}
A.~Krogh ve J.~Hertz. A simple weight decay can improve
generalization. \emph{NeurIPS}, 1991.

\bibitem{srivastava2014dropout}
N.~Srivastava, G.~Hinton, A.~Krizhevsky, I.~Sutskever ve
R.~Salakhutdinov. Dropout: A simple way to prevent neural
networks from overfitting. \emph{JMLR}, 15(1):1929--1958, 2014.

\bibitem{miyato2018spectral}
T.~Miyato, T.~Kataoka, M.~Koyama ve Y.~Yoshida. Spectral
normalization for generative adversarial networks.
\emph{ICLR}, 2018.

\bibitem{vaswani2017attention}
A.~Vaswani ve ark. Attention is all you need. \emph{NeurIPS}, 2017.

\end{thebibliography}

\end{document}
